\documentclass{beamer}
\usetheme{Blubird}

%%%%%%%%%%%%%%%%%%%%%
\usepackage{twemojis}
\newcommand{\lamb}{\texttwemoji{sheep}}

\title{Checking Proofs with Computers}
\subtitle{and the Curry-Howard Correspondence}
\author{Bryan Lu}
\date{July 9, 2024}

  \setcounter{showSlideNumbers}{1}

\begin{document}  

	\setcounter{showProgressBar}{0}
	\setcounter{showSlideNumbers}{0}

\frame{\titlepage}

\begin{frame}{A Little About Me}
  \begin{minipage}{0.6\textwidth}
  \begin{itemize}
    \item HCSSiM 2019 alum
%      \begin{itemize}
%        \item PJ writer/editor
%        \item (former) Frisbee extraordinaire
%      \end{itemize}
    \item Graduate Student in Math @ University of Washington
      \begin{itemize}
        \item BA Math/Computer Science from Cornell University 
      \end{itemize}
    \item Combinatorics! \texttwemoji{blue heart}
    \only<2>{\begin{itemize} 
      \item had a formal verification phase 
  \end{itemize}}
  \end{itemize}  
  \end{minipage}
  \hfill
  \begin{minipage}{0.3\textwidth}
    \includegraphics[width=\textwidth]{profile-pic.png} 
  \end{minipage}
\end{frame}

\section{Context}

  \setcounter{framenumber}{0}
	\setcounter{showProgressBar}{1}
	\setcounter{showSlideNumbers}{1}


\begin{frame}{Computational Power in Math}
  \pause
  \begin{itemize}
    \item The four-color theorem (1976)
    \item Optimal sphere-packing (1998)
    \pause\item ``Happy ending problem'' for hexagon -- 17 points (2006)
    \item Solving any Rubik's Cube -- 20 moves (2010)
    \item Minimum number of clues in sudoku -- 17 (2012)
  \end{itemize} 
\end{frame}


\begin{frame}{Proofs Verified by Computers}
 \small{Proofs/programs checked by \emph{proof assistants}:}
 \pause
 \begin{itemize}
   \item The four-color theorem (Rocq, 2002) 
   \item The CompCert C Compiler (Rocq, 2009)
   \item seL4 kernel (Isabelle/HOL, 2010)
   \item The Feit-Thompson theorem about symsets of odd
     size (Rocq, 2012) \pause 
   \item Peter Scholze, condensed mathematics (Lean, 2021)
   \item Terence Tao, the Polynomial Freiman-Ruzsa conjecture (Lean, 2023)
 \end{itemize}
\end{frame}

\begin{frame}{\includegraphics[width=26pt]{yellowPig.png}}
  \begin{center}
    \Huge{Q: \texttwemoji{thinking} $\Rightarrow$ \texttwemoji{desktop
    computer} $\Rightarrow$ \texttwemoji{check mark
  button}/\texttwemoji{x}?}
\pause 
\vspace{1em}

    \huge{A: \texttwemoji{sparkles}\emph{programming
    languages}\texttwemoji{sparkles}!}
  \end{center}
\end{frame}


\section{Simply-Typed \texorpdfstring{\lamb}{lamb}-da Calculus}

\begin{frame}{Syntax}
Terms in the simply-typed $\lamb$-da calculus:
\[
  e ::= c \mid x \mid \lamb(x:t) \to e \mid e_1 \, e_2 
\] 
Some examples of terms: 
\[
  \ttt{true},\; \ttt{false},\; *, \; x,\; y,\; \lamb(x) \to x, \; (\lamb(x) \to
  \ttt{true}) \; *
\]
\pause
All terms have types:
\[
  t ::= T \mid t_1 \to t_2
\]
Basic types include $\ttt{bool}, \ttt{unit},$ etc. 
\[
  \ttt{true} : \ttt{bool} \quad \ttt{false} : \ttt{bool} \quad * : \ttt{unit}
  \quad \lamb(x) \to \ttt{true} : t \to \ttt{bool}
\] 


\end{frame}


\begin{frame}{Examples: Functions}
  \begin{example}
 \begin{itemize}
   \item $e_1 = \lamb(x:t) \to x \, : \only<1>{\,\; ???} \pause \, t \to t$ \hfill identity function 
    \begin{itemize}
      \item $e_1 \; \ttt{true} = (\lamb (x : t) \to x) \; \ttt{true} = \ttt{true}$  
    \end{itemize}
   \pause \item $e_2 = $ \only<3>{$\lamb(x:t_1) \to \lamb(y:t_2) \to x$}
   \pause {$\lamb (x, y) \to x \, : \only<4>{\,\; ???} \pause  \, \only<5>{t_1 \to (t_2
   \to t_1)} \pause t_1 \to t_2 \to t_1 $} \hfill constant function
   \begin{itemize}
     \item NOT the same as $(t_1 \to t_2) \to t_1$!
   \end{itemize}
  \pause 
\item $e_3 = \lamb (x, y) \to x \; y \, : \only<7>{\,\; ???} \pause \, (t_1 \to t_2) \to t_1 \to
     t_2$ \hfill \footnotesize{evaluation function}
     \begin{itemize}
       \item $e_3 \; e_1 \; * = e_1 \; * = *$
     \end{itemize}
 \end{itemize}
    
  \end{example} 
  
\end{frame}

\begin{frame}{Exercise/Questions?}
  \emph{What is the type of the following term?}
\[
  e_4 = \lamb(x, y, z) \to x \; z \; (y \; z)
\] 
%\pause
%\begin{itemize}
%  \item Suppose $x : \alpha$, $y : \beta$, $z : \gamma$. \pause
%  \item $y$ is a function with one argument of type $\gamma$, so $\beta = \gamma
%    \to \delta$. 
%  \item $x$ is a function with one argument of type $\gamma$ followed by an argument
%    of type $\beta \; \delta = (\gamma \to \delta) \; \gamma = \delta$.
%    \begin{itemize}
%      \item Therefore $\alpha = \gamma \to \delta \to \epsilon$. 
%    \end{itemize} \pause
%  \item $e_4 : (\gamma \to \delta \to \epsilon) \to (\gamma \to \delta) \to
%    \gamma \to \epsilon$. 
%\end{itemize}
%\begin{itemize}
%   \item $\lamb x \, y \, z \to x \, z \, (y \, z)$ -- substitution
%   \item $\lamb x \, y \, z \to (y \, z) \, (x \, z)$ -- contradiction
%   \item $\lamb x \, y \, z \to y \, (x \, z)$ -- contrapositive 
%\end{itemize}
\end{frame}

\begin{frame}{\includegraphics[width=20pt]{yellowPig.png}\texttwemoji{mag_right}}
 \begin{center}
   \includegraphics[width=0.55\textwidth]{more_power.png} 
 \end{center} 
\end{frame}

\section{Let's Talk About Types!}

\begin{frame}{The Natural Numbers}
  Introducing $\NN$ with \emph{inductive types}: 
\begin{center}
\begin{tabular}{l}
  \ttt{inductive} $\NN$ \\
  | $O$ : $\NN$ \\
  | $S$ : $\NN \to \NN$
\end{tabular}
\end{center}
$O$ and $S$ are called \emph{constructors} (also \emph{introduction rules}).  \pause
\vspace{12pt}

\emph{How do these constructors allow us to model $\set{0, 1, 2, \dots}$?} \pause 
\begin{center}
  \begin{tabular}{l}
$O \sim 0 \in \NN$ \qquad $S \sim \text{succ} : \NN \to \NN$ \\ \pause 
$S \; O \sim \text{succ}(0) = 1 \in \NN$ \\ \pause 
$S \; (S \; O) \sim \text{succ}(\text{succ}(0)) = \text{succ}(1) = 2 \in \NN$ \\
\pause 
\vdots
  \end{tabular}
\end{center}
\end{frame}

\begin{frame}{Induction Principles}

  \emph{Induction principles} (\emph{elimination rules}) allow you to define
  functions on inductive types.  
  \pause
 \begin{example}
   The induction principle of $\NN$ requires us to specify values for the ``base case'' and
   ``inductive step.''\pause
  \vspace{12pt}
   \begin{minipage}{0.45\textwidth}
   \begin{tabular}{l}
     \ttt{is\_zero} \, : \, $\NN \to \ttt{bool}$ =\\
     \ttt{ }$\lamb(x : \NN) \to$ \ttt{match}  $x$  \ttt{with} \\
     \ttt{  }| $O \Rightarrow \ttt{true}$ \\
     \ttt{  }| $S \; k \Rightarrow \ttt{false}$ 
   \end{tabular}  
   \end{minipage}\pause
   \begin{minipage}{0.45\textwidth}
   \begin{tabular}{l}
     \ttt{add} \, : \, $\NN \to \NN \to \NN$ =\\
     \ttt{ }$\lamb(x, y) \to$ \ttt{match}  $y$  \ttt{with} \\
     \ttt{  }| $O \Rightarrow x$ \\
     \ttt{  }| $S \; k \Rightarrow S \; (\ttt{add} \; x \; k)$ 
   \end{tabular}  
   \end{minipage}
 \end{example} 
\end{frame}

\begin{frame}{Data Structures}
  Inductive types allow us to model lists of numbers:
  \begin{center}
    \begin{tabular}{l}
        \ttt{inductive natlist}  \\ 
        | \ttt{nil} : \ttt{natlist} \\
        | \ttt{cons} : $\lamb (a : \NN) \to \lamb (t : \ttt{natlist}) \to \ttt{natlist}$ 
    \end{tabular}
  \end{center}\pause
\begin{example} \pause
  \begin{tabular}{rl}
    \ttt{nil} &$\sim$ \quad [ ] \\ \pause 
    \ttt{cons} \, $(S \; O)$ \, \ttt{nil} &$\sim$ \quad [1] \\ \pause
    \ttt{cons} \, $(S \; (S \; O))$ \, (\ttt{cons} \, $(S \; O)$ \, \ttt{nil}) &$\sim$ \pause
    \quad [2, 1] \\
    \pause
                                                             &\vdots
  \end{tabular}
\end{example}
\end{frame}

\begin{frame}{\includegraphics[width=20pt]{yellowPig.png}\texttwemoji{mag_right}}
 \begin{center}
   \includegraphics[width=0.55\textwidth]{more_power.png} 
 \end{center} 
\end{frame}

\begin{frame}{Combined Syntax}
  Combine rules for making terms vs. types into one rule: 
  \[
  t ::= x \mid T \mid \lamb (x :t_1) \to t_2 \mid t_1 \, t_2 \mid \TT
  \] 
  What's changed? \pause 
  \begin{itemize}
    \item Dependent function types
      \begin{itemize}
        \item Functions whose type is $\lamb (x : t_1) \to t_2$, where $t_2$
          may contain $x$
        \item ``$t_1 \to t_2$'' $= \lamb (\_ : t_1) \to t_2$ ($t_2$ doesn't
          depend on $x$)
      \end{itemize}
      \pause
    \item Introducing the \emph{sort} $\TT$, the ``type'' of types
      \begin{itemize}
        \item $\ttt{bool} : \TT$, $\ttt{unit} : \TT$, $\NN : \TT$, etc. 
      \end{itemize}
  \end{itemize}
\end{frame}

\begin{frame}{Application: Lists (again)}
  Recall our type for lists of $\NN$s: 
  \begin{center}
    \begin{tabular}{l}
        \ttt{inductive natlist}  \\ 
        | \ttt{nil} : \ttt{natlist} \\
        | \ttt{cons} : $\lamb (a : \NN) \to \lamb (t : \ttt{natlist}) \to \ttt{natlist}$ 
    \end{tabular}
  \end{center}\pause
  Why not build a list for \emph{any} type? 
  \begin{center}
    \only<2>{
    \begin{tabular}{l}
        \ttt{inductive list} : $\TT \to \TT$ \\ 
        | \ttt{nil} : $\lamb (\alpha : \TT) \to \ttt{list} \; \alpha$ \\
        | \ttt{cons} : $\lamb (\alpha : \TT) \to \lamb (a : \alpha) \to \lamb (t :
        \ttt{list} \; \alpha) \to \ttt{list} \; \alpha$ 
    \end{tabular}}
    \only<3->{
    \begin{tabular}{l}
        \ttt{inductive list} : $\TT \to \TT$ \\ 
        | \ttt{nil} : $\lamb (\alpha : \TT) \to \ttt{list} \; \alpha$ \\
        | \ttt{cons} : $\lamb (\alpha : \TT, a : \alpha, t :
        \ttt{list} \; \alpha) \to \ttt{list} \; \alpha$ 
    \end{tabular}}
  \end{center} 
\end{frame}


\begin{frame}{Terms vs. Types}
 Terms and types can depend on other terms and types: \pause 
  \begin{itemize}
    \item Terms depending on terms: \pause function application \pause 
      \begin{itemize}
        \item $(\lamb(x \; y) \to x) \; w \; z$ depends on $w$. 
      \end{itemize}\pause

    \item Terms depending on types: \pause polymorphic functions \pause
      \begin{itemize}
        \item $(\lamb(x) \to x)$ depends on the type of the input $x$. 
      \end{itemize} \pause

    \item Types depending on types: \pause type constructors \pause
      \begin{itemize}
        \item \ttt{list} $\alpha$ depends on the type $\alpha$. 
      \end{itemize} \pause

    \item Types depending on terms: \emph{dependent types} 
    \begin{itemize}
      \item ??? 
    \end{itemize} 
  \end{itemize}
\end{frame}

\begin{frame}{Example: Type that Depends on a Term}
  Adapt our polymorphic list to record the length of the list!
  Recall:
  \begin{center}
    \begin{tabular}{l}
        \ttt{inductive list} : $\TT \to \TT$ \\ 
        | \ttt{nil} : $\lamb (\alpha : \TT) \to \ttt{list} \; \alpha$ \\
        | \ttt{cons} : $\lamb (\alpha : \TT, a : \alpha, t
        : \ttt{list} \; \alpha) \to \ttt{list} \; \alpha$ 
    \end{tabular}
  \end{center}  
  \pause 
  Now consider:
  \begin{center}
    \begin{tabular}{l}
        \ttt{inductive vector} : $\TT \to \NN \to \TT$ \\ 
        | \ttt{nil} : $\lamb (\alpha : \TT) \to \ttt{vector} \; \alpha \; O$ \\
        | \ttt{cons} : $\lamb (\alpha : \TT, a : \alpha, n : \NN, t : \ttt{vector} \; \alpha \; n) \to \ttt{vector}
        \; \alpha \; (S \; n)$ 
    \end{tabular}
  \end{center} 
  
\end{frame}

\begin{frame}{Examples: Inductive Types}
  Some important inductive dependent types: 
\begin{itemize}
  \item 
   \begin{tabular}{l}
     \ttt{inductive combine} : $\lamb(\alpha : \TT) \to \lamb (\beta : \TT) \to \TT$
 \\
     | \ttt{pair} : $\alpha \to \beta \to \ttt{combine}$ $\alpha$ $\beta$
 \end{tabular}\pause
 \begin{itemize}
   \item \ttt{pair} $(S \; O)$ \ttt{true} : \ttt{combine} $\NN$ \ttt{bool}
 \end{itemize}
\pause
\item \begin{tabular}{l}    
      \ttt{inductive choose} : $\lamb(\alpha : \TT) \to \lamb (\beta : \TT) \to \TT$ \\
    | \ttt{fst} : $\alpha \to \ttt{choose}$ $\alpha$ $\beta$ \\
    | \ttt{snd} : $\beta \to \ttt{choose}$ $\alpha$ $\beta$
\end{tabular}\pause
\begin{itemize}
  \item \ttt{fst} $(S \; (S \; O))$ : $\ttt{choose} \; \NN \; \ttt{bool}$
  \item \ttt{snd false} : $\ttt{choose} \; \NN \; \ttt{bool}$
  \item \ttt{snd false} : $\ttt{choose} \; \ttt{unit bool}$
\end{itemize}\pause

\item \ttt{inductive none} : $\TT$ \pause
  \begin{itemize}
    \item There are NO terms of type \ttt{none}! (\emph{uninhabited})
  \end{itemize}
\end{itemize} 

\end{frame}

\section{Revisiting Logic}

\begin{frame}{Computing with Proofs}
 Suppose $P, Q$ are propositions with proofs $p, q$, respectively. How do we
 computationally get a proof of ...\pause 
 \begin{itemize}
   \item conjunction, $P \wedge Q$? \pause \hfill {\small pair of proofs $(p, q)$.}
     \pause
   \item disjunction, $P \vee Q$? \pause \hfill {\small proof + tag: $(1, p)$ or
   $(2, q)$.}
   \pause
 \item implication, $P \to Q$? \pause \hfill {\small function turning $p \to q$.} 
  \pause
 \item negation, $\neg P$? \pause \hfill {\small function turning $p \to
   \text{(proof of false)}$.} 
 \end{itemize}
\pause
 Suppose $S$ is a set and $T : S \to \text{Prop}$ a proposition over $S$. What
 is a proof of the quantified statement ... \pause
 \begin{itemize}
   \item  $\forall x \in S, T(x)$? \pause \hfill {\small
     function turning $(x \in S) \to \text{proof }t(x)$.}
    \pause
   \item  $\exists x \in S, T(x)$? \pause \hfill {\small
     pair of instance $x \in S$, proof $t(x)$.}
 \end{itemize}

\end{frame}


\section{The Correspondence}

\begin{frame}{The Curry-Howard Correspondence}
  \begin{alertblock}{\Large\centering \emph{types $\rightleftharpoons$ propositions}
    \quad \emph{terms $\rightleftharpoons$ proofs} }
\pause

    \begin{center}
    \begin{tabular}{c|c||c|c}
    Type & Proposition & Expression & Proof  \\ \hline
    \only<2>{???}\only<3->{\ttt{combine} $\alpha$ $\beta$} & $P \wedge Q$ &
    \only<2>{???}\only<3->{$\ttt{pair} \; a \; b$} & $(p, q)$ \\
    \only<2-3>{???}\only<4->{\ttt{choose} $\alpha$ $\beta$} & $P \vee Q$ &
    \only<2-3>{???}\only<4->{$\ttt{fst} \; a, \ttt{snd} \; b$} & $(1, p), (2, q)$\\
    \only<2-4>{???}\only<5->{$\alpha \to \beta$} & $P \to Q$ & 
    \only<2-4>{???}\only<5->{$\lamb : a \mapsto b$} & $f : p \mapsto q$ \\
    \only<2-5>{???}\only<6->{\ttt{unit}} & true & 
    \only<2-5>{???}\only<6->{$*$} & $\exists$ \\
    \only<2-6>{???}\only<7->{\ttt{none}} & false &
    \only<2-6>{???}\only<7->{\texttwemoji{x}} & $\nexists$ \\
    \only<2-7>{???}\only<8->{$\alpha \to \ttt{none}$} & $\neg P$ & 
    \only<2-7>{???}\only<8->{$\lamb : a \mapsto \lightning$} & $f : p
    \mapsto \bot$ \\ 
    \only<2-8>{???}\only<9->{$\lamb(a : \alpha) \to \beta$} & $\forall x \in P$,
    $Q(x)$ &
    \only<2-8>{???}\only<9->{$\lamb : a \mapsto b(a)$} & $f : x \mapsto q(x)$ \\
   ??? & $\exists x \in P$, $Q(x)$ & ??? & $(x, q(x))$
   \end{tabular}
    \end{center}   
 \end{alertblock}
 \begin{center}
  \only<10->{
   \vspace{-1em}
    \Large$\Downarrow$ 


  \large{checking proof correctness $\rightleftharpoons$ type-checking term}}  
 \end{center} 
\end{frame}



\section{Demo!}

\begin{frame}{Resources}
  \begin{itemize}
    \item Rocq
      \begin{itemize}
        \item Software Foundations
      \end{itemize}
    \item Lean 
      \begin{itemize}
        \item The Natural Number Game
        \item leanprover-community (mathlib), Xena Project  
      \end{itemize}
    \item Agda {\scriptsize(less beginner-friendly)}
      \begin{itemize}
        \item HoTTEST Summer School 2022, HoTT Book
        \item agda-unimath, TypeTopology
        \item cubical Agda, 1Lab 
      \end{itemize}
  \end{itemize}
\end{frame}

\section{\texorpdfstring{\includegraphics[width=36pt]{yellowPig.png}}{P.I.G.}}

\backupbegin

\section{Secret Slides}

\begin{frame}{What's the type of a sort?}
  \emph{Can we have $\TT : \TT$?}
  \pause \vspace{12pt}

  \textbf{NO!} It's actually not consistent to allow $\TT :
    \TT$  -- the argument of this
    is similar to ``the barber who shaves everyone who doesn't shave
    themselves'', i.e. Bertrand's paradox.
    Instead, we have \emph{Girard's
    paradox} which breaks this construction. 
    \pause \vspace{12pt}

    In fact, we actually
    need $\TT : \TT_1$, and then $\TT_1 : \TT_2$, etc. The $\TT_i$s are called
    \emph{type universes} and this is more of a technical point that's need to
    make things work, but often isn't relevant.  
\end{frame}


\begin{frame}{The equality type?}
  \emph{What type corresponds to the equality proposition?}
  \pause \vspace{12pt}

  We actually can just make an inductive type for this: 
  \begin{center}
    \begin{tabular}{l}
      \ttt{inductive equal} : $\lamb (\alpha : \TT) \to \alpha \to \alpha \to \TT$ \\
      | \ttt{refl} : $\lamb (\alpha : \TT) \to \lamb(a : \alpha) \to \ttt{equal} \; \alpha \; a \; a $
    \end{tabular}
  \end{center}
  \pause 

  There is only one constructor of this type, and it represents the reflexivity
  of equality (where the two sides of the equality are equivalent). 
\end{frame}


\begin{frame}{Defining complicated propositions?}
  \emph{How do we get more complicated propositions?}
  \pause \vspace{12pt}

  We can just make more types with appropriate constructors. For instance,
  consider the $\leq$ relation on $\NN$:  
  \begin{center}
    \begin{tabular}{l}
      \ttt{inductive leq} : $\lamb (m : \NN) \to \lamb (n : \NN) \to \TT$ \\
      | \ttt{leq-O} : $\lamb (n : \NN) \to \ttt{leq} \; O \; n $ \\
      | \ttt{leq-S} : $\lamb (m : \NN, n : \NN, H:
      \ttt{leq} \; m \; n) \to \ttt{leq} \; (S \; m) \; (S \; n)$
    \end{tabular}
  \end{center}
  \pause 

  Here, we can get a proof of this proposition if we either generate it directly
  from the fact that all naturals are $\geq$ 0, or by reducing successor
  functions on either side of the $\leq$ until we get the above case. This is
  again inductive!
  
\end{frame}

\begin{frame}{Defining algebraic structures?}
  \emph{How does one encode an algebraic structure as a type?}
  \pause \vspace{12pt}

  \emph{Record types} accomplish this, bundling together the
  type for the underlying set and all necessary
  operations/properties. For example, here's a group: 
  \begin{center}
    \begin{tabular}{l}
      \ttt{struct group} : $\lamb (\alpha : \TT) \to \TT$ \quad \{\\
      | \quad \ttt{mul} : $\alpha \to \alpha \to \alpha$ \\
      | \quad \ttt{id} : $\alpha$ \\
      | \quad \ttt{inv} : $\alpha \to \alpha$ \\ 
      | \quad \ttt{mul\_id} : $\lamb (x : \alpha) \to \ttt{mul} \; x \; \ttt{id}
      \; = \; x$ $\dots$ \} \\
    \end{tabular}
  \end{center}
  \pause

  One can also encode hierarchies of algebraic structures in types as well,
  e.g. a ring is a group + structure. Record types are also useful for
  encoding $\exists$ propositions. 
\end{frame}

\backupend











\end{document}

